% ===== Topic 2, part b: problems 16-30 (EN) =====
% helper: 2x2 grid with one letter; #1,#2 = centre of the cell, #3 = letter, #4 = node transform
\providecommand{\tbGrid}[4]{\begin{tikzpicture}[baseline=0.2cm,x=0.5cm,y=0.5cm]
  \draw[Ink,line width=0.7pt] (0,0) grid[step=1] (2,2);
  \node[font=\sffamily\bfseries,text=tmBlue!80!black,inner sep=0pt,#4] at (#1,#2) {#3};\end{tikzpicture}}

% ---------------------------------------------------------------- 16
\begin{probbox}{16}{\relax}
A rectangular garden has a length that is $6$ m more than \emph{twice} its width.
The perimeter of the garden is $120$ m. How many metres wide is the garden?
\choices{$8$ m}{$18$ m}{$27$ m}{$38$ m}{$22$ m}
\end{probbox}
\probsol{
\textbf{Understanding the problem.} We know two things: (1) length $=$ twice the width, plus $6$ m;
(2) perimeter $=120$ m. We must find the width. We introduce an unknown and form an equation.

\step{1}{Choose the unknown and express the length in terms of it}
Let the width of the garden be $x$ (m), with $x>0$.
``Twice the width'' is $2x$; the length is $6$ m more than that, so the length is $2x+6$ (m).
\begin{figC}
\draw[Ink,thick,fill=tmGreen!10] (0,0) rectangle (4.2,1.8);
\node[below] at (2.1,0){$2x+6$};
\node[right] at (4.2,0.9){$x$};
\end{figC}

\begin{recallbox}Perimeter of a rectangle $=(\text{length}+\text{width})\cdot 2$.\end{recallbox}

\step{2}{Form the equation}
The perimeter is $120$ m, so
\[ \big((2x+6)+x\big)\cdot 2=120. \]

\step{3}{Solve the equation}
Divide both sides by $2$: $(2x+6)+x=60$.
Collect like terms: $3x+6=60$.
Subtract $6$ from both sides: $3x=60-6=54$.
Divide both sides by $3$: $x=54\div 3=18$ (this satisfies $x>0$).

So the width is $18$ m and the length is $2\cdot 18+6=42$ (m).

\begin{warnbox}Many students forget that the perimeter contains \emph{two} lengths and \emph{two} widths, write
$3x+6=120$ and get $x=38$ -- exactly the ``trap'' option (D).\end{warnbox}

\verify{Length $42$ m, width $18$ m: $42=2\cdot 18+6$ (correct) and the perimeter is
$(42+18)\cdot 2=60\cdot 2=120$ m (correct).}
A width of $18$ m is option \textbf{(B)}.
\ans{(B) $18$ m}
}

% ---------------------------------------------------------------- 17
\begin{probbox}{17}{\relax}
The letter P is written in a $2\times 2$ grid of squares as shown: \tbGrid{1.5}{1.5}{P}{}\,.
A sequence of rotations about the centre of the grid and reflections in the two lines through the centre
of the grid (the horizontal line and the vertical line) is carried out, and the result is:
\tbGrid{0.5}{0.5}{P}{cm={0,-1,-1,0,(0,0)}}\,.

If \emph{exactly the same sequence of transformations} is applied to the grid \tbGrid{0.5}{1.5}{A}{}, which
picture do we get?
\par\smallskip\noindent
\begin{tabularx}{\linewidth}{@{}XXXXX@{}}
(A) \tbGrid{0.5}{1.5}{A}{} & (B) \tbGrid{0.5}{1.5}{A}{rotate=90} & (C) \tbGrid{1.5}{0.5}{A}{rotate=-90} &
(D) \tbGrid{1.5}{0.5}{A}{} & (E) \tbGrid{1.5}{0.5}{A}{rotate=180}
\end{tabularx}
\end{probbox}
\probsol{
\textbf{Understanding the problem.} We do not know which individual steps the sequence contains, but we know
its \emph{final effect} on the letter P. Idea: imagine the grid is a square \emph{card} made of clear plastic
with the letter drawn on it. Every rotation or reflection just turns or flips the card and puts it back in the
same place. The whole sequence is therefore also just \emph{one way of putting the card back}. We read off that
way from the picture of P, and then apply it to A.

\step{1}{See where the letter P went}
The letter P in the \emph{top-right} cell has moved to the \emph{bottom-left} cell. At the same time, the
straight stroke of P (its ``stem''), which was vertical, is now horizontal at the top; the round ``bowl'', which
was to the right of the stem, is now underneath it.

\step{2}{Find out how the card was put back}
There are two ways of putting the card back that move the top-right cell to the bottom-left cell:
\begin{itemize}[leftmargin=1.4em,itemsep=1pt]
\item a $180^\circ$ rotation: then P would be upside down (its stem still \emph{vertical}, looking like a ``d'');
\item a reflection in the diagonal joining the top-left corner to the bottom-right corner (the dashed line in the
picture below): then the vertical stem of P becomes \emph{horizontal}.
\end{itemize}
In the given picture the stem is horizontal, so the sequence acts like the
\textbf{reflection in the top-left -- bottom-right diagonal}.
\begin{figC}\begin{scope}[x=0.8cm,y=0.8cm]
\draw[Ink,line width=0.8pt] (0,0) grid[step=1] (2,2);
\draw[tmOrange,dashed,line width=1pt] (-0.3,2.3)--(2.3,-0.3);
\node[font=\sffamily\bfseries\large,text=tmBlue] at (1.5,1.5){P};
\node[font=\sffamily\bfseries\large,text=tmGreen!80!black,cm={0,-1,-1,0,(0,0)}] at (0.5,0.5){P};
\draw[->,tmPurple,thick] (1.55,1.15) to[bend left=35] (0.85,0.45);
\node[font=\footnotesize,text=tmOrange!80!black,anchor=west] at (2.3,-0.3){mirror line};
\begin{scope}[xshift=4.2cm]
\draw[Ink,line width=0.8pt] (0,0) grid[step=1] (2,2);
\draw[tmOrange,dashed,line width=1pt] (-0.3,2.3)--(2.3,-0.3);
\node[font=\sffamily\bfseries\large,text=tmBlue] at (0.5,1.5){A};
\node[font=\footnotesize,text=Slate] at (1,-0.45){A lies on the line};
\end{scope}
\end{scope}\end{figC}
\begin{recallbox}A sequence of rotations about the centre and reflections in the horizontal and vertical lines can
have the same effect as a reflection in a diagonal. For example: \emph{a $90^\circ$ clockwise rotation followed
by a reflection in the vertical line}. Check with P: the $90^\circ$ rotation takes P from the top-right cell down
to the bottom-right cell, the reflection in the vertical line takes it to the bottom-left cell, and the stem of P
is horizontal -- exactly as in the given picture.\end{recallbox}

\step{3}{Apply it to the letter A}
The top-left cell lies \emph{right on} the mirror line (the diagonal passes through it), so the letter A
\emph{stays in the top-left cell}. This diagonal reflection turns the direction ``up'' into the direction
``left''. The tip of A, which pointed up, now points left. The letter A is symmetric left-to-right, so we need
not worry about it being ``mirrored''.

\begin{warnbox}Do not look only at the position of the cell! P moved from the top-right to the bottom-left
cell, but that does \emph{not} mean A also ``moves diagonally down''. First find the transformation of
\emph{the whole card}, then apply it.\end{warnbox}

\verify{Use the concrete sequence ``$90^\circ$ clockwise rotation, then reflection in the vertical line'': the
rotation moves A from the top-left to the top-right cell with its tip pointing right; the reflection in the
vertical line brings A back to the top-left cell with its tip pointing left. This is picture (B).}
\ans{(B) A in the top-left cell, tip pointing left}
}

% ---------------------------------------------------------------- 18
\begin{probbox}{18}{\relax}
In a square with side $8$ cm, the two diagonals are drawn. What is the area of the shaded region?
\begin{figC}\begin{scope}[x=0.38cm,y=0.38cm]
\fill[tmBlue!35] (0,0)--(8,0)--(4,4)--cycle;
\draw[Ink,line width=1pt] (0,0) rectangle (8,8);
\draw[Ink,line width=0.9pt] (0,0)--(8,8) (0,8)--(8,0);
\node[left] at (0,4){$8$ cm};
\end{scope}\end{figC}
\choices{$4$ cm$^2$}{$8$ cm$^2$}{$16$ cm$^2$}{$56$ cm$^2$}{$64$ cm$^2$}
\end{probbox}
\probsol{
\textbf{Understanding the problem.} The two diagonals divide the square into $4$ triangles. The shaded region is
one of these four triangles. We need its area.

\step{1}{Compute the area of the square}
\begin{recallbox}The area of a square with side $a$ is $a\times a$.\end{recallbox}
$A_{\text{square}}=8\times 8=64$ (cm$^2$).

\step{2}{Show that the four triangles have equal areas}
\begin{proofbox}The two diagonals of a square are equal and bisect each other at a point $O$, so $O$ is the same
distance from all four vertices. Rotating the square $90^\circ$ about $O$ moves each triangle exactly onto the
next one. So the four triangles are identical, and each is $\tfrac14$ of the square.\end{proofbox}

\step{3}{Compute the shaded area}
$A=64\div 4=16$ (cm$^2$).

\verify{Direct calculation: the shaded triangle has the bottom side as its base, $8$ cm, and its height is the
distance from the centre $O$ to that side, which is half the side, $4$ cm. Area $=\dfrac{8\times 4}{2}=16$
cm$^2$. It matches.}
Option \textbf{(C)}.
\ans{(C) $16$ cm$^2$}
}

% ---------------------------------------------------------------- 19
\begin{probbox}{19}{\relax}
The grid below is made of identical small squares. What percentage of the area of the whole grid is shaded?
\begin{figC}\begin{scope}[x=0.36cm,y=0.36cm]
\foreach \r in {0,...,7}{\foreach \c in {0,...,7}{\fill[tmBlue!40] (\c,\r) rectangle (\c+1,\r+1);}}
\foreach \c/\r in {2/7,5/7,2/6,5/6,0/5,1/5,6/5,7/5,0/2,1/2,6/2,7/2,2/1,5/1,2/0,5/0}{\fill[white] (\c,\r) rectangle (\c+1,\r+1);}
\draw[Ink,line width=0.7pt] (0,0) grid[step=1] (8,8);
\end{scope}\end{figC}
\choices{$84\%$}{$80\%$}{$90\%$}{$75\%$}{$66\%$}
\end{probbox}
\probsol{
\textbf{Understanding the problem.} To find a percentage we need to know how many squares there are in total and
how many are shaded. Counting the \emph{white} squares is quicker because there are fewer of them.

\step{1}{Count all the squares}
The grid has $8$ rows of $8$ squares: $8\times 8=64$ squares.

\step{2}{Count the white squares}
The white squares form $8$ ``pairs'' of two adjacent squares, placed symmetrically on the four sides:
\begin{itemize}[leftmargin=1.4em,itemsep=1pt]
\item top two rows: columns $3$ and $6$ have a white square in both rows: $4$ squares;
\item bottom two rows: likewise, $4$ squares;
\item row $3$ (from the top): $2$ white squares at the left edge and $2$ at the right edge: $4$ squares;
\item row $6$: likewise, $4$ squares.
\end{itemize}
Total number of white squares: $4+4+4+4=16$.

\step{3}{Count the shaded squares and convert to a percentage}
Number of shaded squares: $64-16=48$.
\begin{recallbox}The percentage that $a$ is of $b$ equals $\dfrac{a}{b}\times 100\%$.\end{recallbox}
\[\frac{48}{64}=\frac{3}{4}=\frac{75}{100}=75\%.\]

\begin{tipbox}You do not even need to count everything: the white squares make up $\dfrac{16}{64}=\dfrac14=25\%$,
so the shaded squares make up $100\%-25\%=75\%$.\end{tipbox}

\verify{$75\%$ of $64$ is $64\times\dfrac34=48$ squares, and $48+16=64$. Correct.}
Option \textbf{(D)}.
\ans{(D) $75\%$}
}

% ---------------------------------------------------------------- 20
\begin{probbox}{20}{\relax}
$JKLM$ is a square and $PQRS$ is a rectangle; they overlap as shown in the diagram. Given that $JK$ is parallel
to $PQ$, $JK=8$ and $PS=2$, find the total area of the shaded regions.
\begin{figC}\begin{scope}[x=0.3cm,y=0.3cm]
\fill[tmBlue!35] (5,-3) rectangle (13,5);
\draw[Ink,line width=0.9pt] (5,-3) rectangle (13,5);
\fill[white] (0,0) rectangle (18,2);
\draw[Ink,line width=0.9pt] (0,0) rectangle (18,2);
\node[above left] at (0,2){$P$};\node[above right] at (18,2){$Q$};
\node[below right] at (18,0){$R$};\node[below left] at (0,0){$S$};
\node[above] at (5,5){$J$};\node[above] at (13,5){$K$};
\node[below] at (13,-3){$L$};\node[below] at (5,-3){$M$};
\end{scope}\end{figC}
\choices{$32$}{$16$}{$56$}{$48$}{$62$}
\end{probbox}
\probsol{
\textbf{Understanding the problem.} The shaded part is the part of the square that is \emph{not} covered by the
rectangular strip. So: shaded area $=$ area of the square $-$ area of the overlap of the two figures.

\step{1}{Area of the square $JKLM$}
\begin{recallbox}Area of a square $=$ side $\times$ side; area of a rectangle $=$ length $\times$ width.\end{recallbox}
$A_{JKLM}=8\times 8=64$.

\step{2}{Identify the overlap}
Since $JK\parallel PQ$, the sides of the strip $PQRS$ are parallel to the side $JK$ of the square, i.e.\ the strip
runs ``horizontally'' across the square. The strip is longer than the square, so it crosses the whole width of
the square. The overlap is a rectangle with:
\begin{itemize}[leftmargin=1.4em,itemsep=1pt]
\item length equal to the side of the square: $8$;
\item width equal to the width of the strip: $PS=2$.
\end{itemize}
Area of the overlap: $8\times 2=16$.

\step{3}{Compute the shaded area}
$64-16=48$.

\begin{warnbox}We do not need to know how high or low the strip is, or how big each shaded piece is: their total
is always $64-16=48$. Many students only compute $8\times 2=16$ (option B) -- that is the \emph{covered} part, not
the shaded part.\end{warnbox}

\verify{Suppose the strip is $3$ below the top side $JK$; the top piece has area $8\times 3=24$, the bottom piece
has height $8-3-2=3$, area $24$. Total $48$. Correct.}
Option \textbf{(D)}.
\ans{(D) $48$}
}

% ---------------------------------------------------------------- 21
\begin{probbox}{21}{\relax}
A triangle can be built with side lengths $4$, $5$, $8$, but \emph{no} triangle can be built with side lengths
$4$, $5$, $10$. Using only the lengths $2$, $3$, $5$, $7$, $11$ (each length may be used more than once), how
many different triangles with \emph{exactly two equal sides} can be built?
\choices{$8$}{$5$}{$20$}{$10$}{$14$}
\end{probbox}
\probsol{
\textbf{Understanding the problem.} A triangle with exactly two equal sides is isosceles but not equilateral: two
legs of length $a$ and a base of length $b$ with $b\neq a$. We must count the pairs $(a,b)$ taken from
$\{2,3,5,7,11\}$ such that the three segments $a,a,b$ really form a triangle.

\begin{recallbox}\textbf{Triangle inequality.} Three segments are the sides of a triangle exactly when the sum of
any two of them is greater than the third. In practice it is enough to check that
\emph{the two shorter segments add up to more than the longest one}. Example: $4+5=9>8$, so it works;
$4+5=9<10$, so it does not.\end{recallbox}

\step{1}{Simplify the condition for an isosceles triangle $a,a,b$}
The sums to check are $a+a>b$ and $a+b>a$. The second is always true (since $b>0$). So we only need
\[ 2a>b \quad\text{(twice the leg is greater than the base).}\]

\step{2}{Make a table by the leg $a$}
\begin{center}
\begin{tabular}{|c|c|c|c|}
\hline
Leg $a$ & $2a$ & Base $b\ne a$ with $b<2a$ & Number of triangles\\ \hline
$2$ & $4$ & $3$ & $1$\\
$3$ & $6$ & $2,\ 5$ & $2$\\
$5$ & $10$ & $2,\ 3,\ 7$ & $3$\\
$7$ & $14$ & $2,\ 3,\ 5,\ 11$ & $4$\\
$11$ & $22$ & $2,\ 3,\ 5,\ 7$ & $4$\\ \hline
\end{tabular}
\end{center}
For instance, with $a=3$: the base $b=7$ does not work because $3+3=6<7$; $b=11$ does not work either.
With $a=5$: $b=11$ does not work because $5+5=10<11$.

\step{3}{Add up}
$1+2+3+4+4=14$ triangles.

\begin{warnbox}Triangles such as $2,2,5$ or $3,3,7$ ``look'' possible, but in fact the two legs are too short to
meet. Also, the triangle $5,5,5$ has \emph{three} equal sides, so it is not counted.\end{warnbox}

\verify{Count the other way: there are $5\times 4=20$ pairs $(a,b)$ with $b\neq a$. The rejected pairs (because
$b\ge 2a$) are $(2,5),(2,7),(2,11),(3,7),(3,11),(5,11)$: $6$ pairs. $20-6=14$. It matches.}
Option \textbf{(E)}.
\ans{(E) $14$}
}

% ---------------------------------------------------------------- 22
\begin{probbox}{22}{\relax}
Parallelogram $PQRS$ has area $40$. Points $T$ and $V$ are the midpoints of sides $PS$ and $RS$ respectively.
Find the area of quadrilateral $PRVT$.
\begin{figC}\begin{scope}[x=0.55cm,y=0.55cm]
\fill[tmBlue!25] (0,0)--(8,4)--(7,2)--(3,0)--cycle;
\draw[Ink,line width=1pt] (0,0)--(2,4)--(8,4)--(6,0)--cycle;
\draw[Ink,dashed,line width=0.8pt] (0,0)--(8,4) (3,0)--(7,2);
\foreach \x in {1.5,4.5}{\draw[Ink] (\x,-0.15)--(\x,0.15);}
\foreach \x/\y in {6.5/1,7.5/3}{\draw[Ink] (\x-0.15,\y+0.07)--(\x+0.15,\y-0.07);\draw[Ink] (\x-0.15,\y+0.17)--(\x+0.15,\y+0.03);}
\node[below left] at (0,0){$P$};\node[above left] at (2,4){$Q$};
\node[above right] at (8,4){$R$};\node[below right] at (6,0){$S$};
\node[below] at (3,0){$T$};\node[right] at (7,2){$V$};
\end{scope}\end{figC}
\choices{$10$}{$12$}{$15$}{$16$}{$18$}
\end{probbox}
\probsol{
\textbf{Understanding the problem.} Quadrilateral $PRVT$ is what remains when the small triangle $TVS$ at
corner $S$ is removed from triangle $PRS$. So we compute $A_{PRS}$ and $A_{TVS}$.

\step{1}{Area of triangle $PRS$}
\begin{recallbox}A diagonal divides a parallelogram into two identical triangles, so each triangle has half the
area of the parallelogram.\end{recallbox}
$A_{PRS}=40\div 2=20$.

\step{2}{Area of triangle $TRS$ (join $T$ to $R$)}
\begin{recallbox}Two triangles with \emph{the same height} have areas in the same ratio as their bases. (Because
the area of a triangle $=\tfrac12\times$ base $\times$ height.)\end{recallbox}
Triangles $TRS$ and $PRS$ share the vertex $R$, and their bases $TS$ and $PS$ lie on the same line, so they have
the same height from $R$. Since $T$ is the midpoint of $PS$, $TS=\tfrac12 PS$. Therefore
$A_{TRS}=\tfrac12\cdot 20=10$.

\step{3}{Area of triangle $TVS$}
Triangles $TVS$ and $TRS$ share the vertex $T$, and their bases $VS$ and $RS$ lie on the line $RS$.
Since $V$ is the midpoint of $RS$, $VS=\tfrac12 RS$, so $A_{TVS}=\tfrac12\cdot 10=5$.

\step{4}{Area of quadrilateral $PRVT$}
$A_{PRVT}=A_{PRS}-A_{TVS}=20-5=15$.

\begin{tipbox}A ``half -- half'' triangle at a corner has area $\tfrac12\times\tfrac12=\tfrac14$ of the big
triangle. This is a very useful fact!\end{tipbox}

\verify{Place the figure on a grid: $P(0,0)$, $S(6,0)$, $R(8,4)$, $Q(2,4)$ (area $6\times4=24$, not $40$, but the
ratios are the same). Then $T(3,0)$, $V(7,2)$; $A_{PRS}=12$, $A_{TVS}=\tfrac12\cdot 3\cdot 2=3$, so
$A_{PRVT}=9=\tfrac{3}{8}\cdot 24$. For the given figure: $\tfrac38\cdot 40=15$. It matches.}
Option \textbf{(C)}.
\ans{(C) $15$}
}

% ---------------------------------------------------------------- 23
\begin{probbox}{23}{\stars{1}}
On a $10\times 10$ grid of squares (each square has side $1$), triangle $FGH$ is drawn as shown. Point $P$ can be
any of the $41$ points where gridlines cross \emph{inside} (not on the sides of) triangle $FGH$. For each position
of $P$ we get exactly three triangles: $FPG$, $GPH$, $HPF$. Of all these $123$ triangles, how many have an area
\emph{exactly half} the area of triangle $FGH$?
\begin{figC}\begin{scope}[x=0.36cm,y=0.36cm]
\draw[Line,line width=0.5pt] (0,0) grid[step=1] (10,10);
\draw[Ink,line width=1.3pt] (0,0)--(10,0)--(5,10)--cycle;
\node[above] at (5,10){$F$};\node[below left] at (0,0){$G$};\node[below right] at (10,0){$H$};
\end{scope}\end{figC}
\choices{$5$}{$7$}{$3$}{$11$}{$9$}
\end{probbox}
\probsol{
\textbf{Understanding the problem.} Use coordinates: $G(0,0)$, $H(10,0)$, $F(5,10)$. Triangle $FGH$ has base
$GH=10$ and height $10$, so
\[ A_{FGH}=\frac{10\times 10}{2}=50,\qquad \text{and half of it is } 25. \]
We count, for each type of triangle $GPH$, $FPG$, $HPF$ in turn, the positions of $P$ that give area $25$.

\begin{recallbox}Two triangles with \emph{the same base} have areas in the same ratio as their heights.
Points that are the same distance from a line (on the same side) lie on a line parallel to it.\end{recallbox}

\step{1}{Triangle $GPH$}
Triangle $GPH$ has the same base $GH$ as triangle $FGH$. For half the area, the height from $P$ must be half of
the height $10$, i.e.\ $5$. So $P$ lies on the 5th horizontal gridline (counting from $GH$).
On this line, side $FG$ passes through the point with $x$-coordinate $2.5$ and side $FH$ through the point with
$x$-coordinate $7.5$ (the midpoints of those sides). The grid points \emph{inside} have $x=3,4,5,6,7$:
\[ (3,5),\ (4,5),\ (5,5),\ (6,5),\ (7,5)\quad\Rightarrow\quad 5 \text{ triangles.} \]

\step{2}{Triangle $FPG$}
This time the common base is $FG$. Let $M(5,0)$ be the midpoint of $GH$ and $N(7.5,5)$ the midpoint of $FH$.
\begin{proofbox}Triangles $FMG$ and $FHG$ have the same height from $F$, and the base $GM=\tfrac12 GH$, so
$A_{FMG}=\tfrac12 A_{FHG}=25$. So point $M$ is exactly half as far from $FG$ as $H$ is. The same holds for $N$
(the midpoint of $FH$). Hence $A_{FPG}=25$ exactly when $P$ is the same distance from $FG$ as $M$, i.e.\ when $P$
lies on the line $MN$ (parallel to $FG$).\end{proofbox}
Going from $G$ to $F$ we move $5$ right and $10$ up, i.e.\ ``$1$ right, $2$ up''. Line $MN$ is parallel to $FG$,
so from $M(5,0)$ we also go ``$1$ right, $2$ up'': we meet $(6,2)$, $(7,4)$, and then $N(7.5,5)$ (not a grid
point). The two points $(6,2)$, $(7,4)$ are inside the triangle $\Rightarrow$ $2$ triangles.

\step{3}{Triangle $HPF$}
The figure is symmetric about the vertical line $x=5$, so the result is the same as in step 2: $P=(4,2)$ or
$(3,4)$ $\Rightarrow$ $2$ triangles.

\step{4}{Add up}
$5+2+2=9$ triangles.

\verify{Check $P(6,2)$ with the ``enclosing rectangle'' method: the rectangle from $x=0$ to $6$ and $y=0$ to $10$
has area $60$. Remove the three right triangles outside triangle $FPG$: $\tfrac{5\cdot 10}{2}=25$,
$\tfrac{1\cdot 8}{2}=4$, $\tfrac{6\cdot 2}{2}=6$. What remains is $60-35=25$, exactly half of $50$.}
\begin{warnbox}Do not forget the two types $FPG$ and $HPF$! If you only consider triangles with the horizontal base
$GH$, you will wrongly choose option (A) $5$.\end{warnbox}
Option \textbf{(E)}.
\ans{(E) $9$}
}

% ---------------------------------------------------------------- 24
\begin{probbox}{24}{\stars{1}}
Points $T,U,V,W,X,Y$ lie on the sides of square $PQRS$ as shown, with
\[PT=TU=UQ=QV=VW=WR=XS=SY.\]
What fraction of the area of square $PQRS$ is shaded?
\begin{figC}\begin{scope}[x=1.25cm,y=1.25cm]
\fill[tmBlue!35] (1,3)--(2,3)--(1.5,2.5)--(2.5,1.5)--(3,2)--(3,1)--(1,1)--cycle;
\draw[Ink,line width=1pt] (0,0) rectangle (3,3);
\draw[Ink,line width=0.9pt] (1,3)--(1,0) (0,1)--(2,3) (1,3)--(3,1) (1,0)--(3,2) (0,1)--(3,1);
\foreach \x in {0.5,1.5,2.5}{\draw[Ink] (\x-0.03,2.92)--(\x-0.03,3.08) (\x+0.03,2.92)--(\x+0.03,3.08);}
\foreach \y in {0.5,1.5,2.5}{\draw[Ink] (2.92,\y-0.03)--(3.08,\y-0.03) (2.92,\y+0.03)--(3.08,\y+0.03);}
\draw[Ink] (0.47,-0.08)--(0.47,0.08) (0.53,-0.08)--(0.53,0.08);
\draw[Ink] (-0.08,0.47)--(0.08,0.47) (-0.08,0.53)--(0.08,0.53);
\node[above left] at (0,3){$P$};\node[above right] at (3,3){$Q$};\node[below right] at (3,0){$R$};
\node[below left] at (0,0){$S$};\node[above] at (1,3){$T$};\node[above] at (2,3){$U$};
\node[right] at (3,2){$V$};\node[right] at (3,1){$W$};\node[below] at (1,0){$X$};\node[left] at (0,1){$Y$};
\end{scope}\end{figC}
\par\smallskip\noindent
\begin{tabularx}{\linewidth}{@{}XXXXX@{}}(A)\ $\dfrac{5}{18}$&(B)\ $\dfrac13$&(C)\ $\dfrac29$&(D)\ $\dfrac14$&(E)\ $\dfrac16$\end{tabularx}
\end{probbox}
\probsol{
\textbf{Understanding the problem.} Each side of the square is divided into $3$ equal parts. Let each part have
length $1$ unit; then the square has side $3$ and area $3\times 3=9$. A nice idea: draw extra lines dividing the
square into a $3\times 3$ grid of unit squares (dotted lines); then every line in the figure becomes very ``tidy''.
\begin{figC}\begin{scope}[x=1.25cm,y=1.25cm]
\fill[tmBlue!12] (1,3)--(3,3)--(3,1)--(1,1)--cycle;
\fill[tmBlue!40] (1,3)--(3,1)--(1,1)--cycle;
\fill[tmOrange!50] (1,3)--(2,3)--(1.5,2.5)--cycle;
\fill[tmGreen!50] (3,2)--(3,1)--(2.5,1.5)--cycle;
\draw[Slate,densely dotted] (0,0) grid[step=1] (3,3);
\draw[Ink,line width=1pt] (0,0) rectangle (3,3);
\draw[Ink,line width=0.9pt] (1,3)--(1,0) (0,1)--(2,3) (1,3)--(3,1) (1,0)--(3,2) (0,1)--(3,1);
\draw[tmOrange!80!black,dashed] (1,2)--(2,3) (1,3)--(2,2);
\draw[tmGreen!70!black,dashed] (2,2)--(3,1) (2,1)--(3,2);
\node[above left] at (0,3){$P$};\node[above right] at (3,3){$Q$};\node[below right] at (3,0){$R$};
\node[below left] at (0,0){$S$};\node[above] at (1,3){$T$};\node[above] at (2,3){$U$};
\node[right] at (3,2){$V$};\node[right] at (3,1){$W$};\node[below] at (1,0){$X$};\node[left] at (0,1){$Y$};
\node[below left] at (1,1){$Z$};\node[font=\footnotesize] at (1.6,1.45){$2$};
\end{scope}\end{figC}

\step{1}{Recognise the three pieces of the shaded region}
Let $Z$ be the point where $TX$ meets $YW$. The shaded region consists of three non-overlapping pieces: the big
triangle $TZW$ (dark blue), the small triangle next to side $TU$ (orange) and the small triangle next to side
$VW$ (green).

\step{2}{Triangle $TZW$}
$TQWZ$ is a $2\times 2$ square (made of $4$ unit squares) and $TW$ is its diagonal.
\begin{recallbox}A diagonal divides a square (rectangle) into two triangles of equal area.\end{recallbox}
So $A_{TZW}=\tfrac12\cdot 4=2$.

\step{3}{The two small triangles}
Look at the unit square whose top side is $TU$. Line $YU$ goes ``$1$ right, $1$ up'' through grid corners, so it
contains one diagonal of this unit square; line $TW$ goes ``$1$ right, $1$ down'' from $T$, so it contains the
other diagonal. The two diagonals of a square divide it into $4$ equal triangles, and the orange triangle is one
of them: area $\tfrac14\cdot 1=\tfrac14$. In exactly the same way, the green triangle lies in the unit square whose
right side is $VW$, cut by the two diagonals $XV$ and $TW$: its area is also $\tfrac14$.

\step{4}{Compute the fraction}
Shaded area: $2+\tfrac14+\tfrac14=\tfrac52$. The required fraction:
\[ \frac{5}{2}\div 9=\frac{5}{2}\cdot\frac{1}{9}=\frac{5}{18}. \]

\begin{warnbox}Do not ``estimate by eye''! The shaded part looks close to $\tfrac13$, but the exact calculation
gives $\tfrac{5}{18}$ (a little less than $\tfrac{6}{18}=\tfrac13$).\end{warnbox}
\verify{Put $S(0,0)$ and side $3$; computing the area of the shaded polygon by adding the triangles gives the same
result $\tfrac52$, and $\tfrac{5}{2}=\tfrac{5}{18}\cdot 9$. Correct.}
Option \textbf{(A)}.
\ans{(A) $\dfrac{5}{18}$}
}

% ---------------------------------------------------------------- 25
\begin{probbox}{25}{\relax}
Rectangles measuring $4\times 2$ are arranged by the following rule: from the second rectangle on, the top-left
vertex of each rectangle is placed exactly at the midpoint of the bottom side of the rectangle just above it (as
shown). When $10$ rectangles in total are arranged in this way, what is the perimeter of the whole figure?
\begin{figC}\begin{scope}[x=0.42cm,y=0.42cm]
\foreach \i in {0,1,2,3}{\draw[Ink,line width=0.9pt] (2*\i,-2*\i) rectangle (2*\i+4,-2*\i+2);}
\foreach \i in {0,1,2}{\draw[Ink] (2*\i+1,-2*\i-0.25)--(2*\i+1,-2*\i+0.25) (2*\i+3,-2*\i-0.25)--(2*\i+3,-2*\i+0.25);}
\node[above] at (2,2){$4$};\node[left] at (0,1){$2$};
\foreach \k in {0,1,2}{\fill[Ink] (9+0.6*\k,-6.8-0.6*\k) circle (2.5pt);}
\end{scope}\end{figC}
\choices{$48$}{$64$}{$90$}{$84$}{$100$}
\end{probbox}
\probsol{
\textbf{Understanding the problem.} Each new rectangle is shifted $2$ to the right (half the length) and $2$ down
(exactly the width). The whole figure looks like a ``staircase''. We compute the perimeter in two ways.

\step{1}{Method 1: count each kind of side}
\emph{Vertical sides.} Each rectangle has a left side and a right side of length $2$, and both lie on the
boundary (not covered by another rectangle). Total: $10\times(2+2)=40$.

\emph{Horizontal sides on top.} The first rectangle shows its whole top side: $4$. For each later rectangle, the
left half of its top side is against the bottom side of the previous one, so only the right half shows: $2$.
Total: $4+9\times 2=22$.

\emph{Horizontal sides at the bottom.} Each of the first $9$ rectangles shows only the left half of its bottom
side: $2$; the last one shows its whole bottom side: $4$. Total: $9\times 2+4=22$.

Perimeter: $40+22+22=84$.

\step{2}{Method 2: ``straighten'' the staircase}
\begin{proofbox}Walking along the top-right boundary from the top-left corner of the first rectangle to the
bottom-right corner of the last one, we only ever go right or down. The horizontal pieces add up to exactly the
width of the whole figure, and the vertical pieces add up to exactly its height. The bottom-left boundary is the
same. Hence the perimeter of the staircase \emph{equals the perimeter of the smallest rectangle enclosing
it}.\end{proofbox}
Width: $4+9\times 2=22$; height: $10\times 2=20$. Perimeter: $(22+20)\times 2=84$.

\begin{warnbox}Do not take $10\times$ the perimeter of one rectangle ($10\times 12=120$): the pieces where two
rectangles touch are \emph{inside}, not on the boundary.\end{warnbox}
\verify{Try $2$ rectangles: method 1 gives $8+(4+2)+(2+4)=20$; the enclosing rectangle is $6\times 4$, perimeter
$20$. With $n$ rectangles the perimeter is $2\big((4+2(n-1))+2n\big)=8n+4$; $n=10$ gives $84$. It matches.}
Option \textbf{(D)}.
\ans{(D) $84$}
}

% ---------------------------------------------------------------- 26
\begin{probbox}{26}{\relax}
A quadrilateral has vertices $A(7,6)$, $B(-5,1)$, $C(-2,-3)$ and $D(10,2)$ in the coordinate plane.
The area of the quadrilateral (in square units) is
\choices{$60$}{$63$}{$67$}{$70$}{$72$}
\end{probbox}
\probsol{
\textbf{Understanding the problem.} The quadrilateral has ``slanted'' sides, and there is no ready-made formula.
A standard technique: \emph{enclose the quadrilateral in a rectangle with sides parallel to the axes}, then
subtract the right triangles in the four corners.
\begin{figC}\begin{scope}[x=0.4cm,y=0.4cm]
\fill[tmOrange!15] (-5,-3) rectangle (10,6);
\fill[tmBlue!30] (7,6)--(-5,1)--(-2,-3)--(10,2)--cycle;
\draw[Line] (-5,-3) grid[step=1] (10,6);
\draw[->,Slate] (-6,0)--(11,0) node[right]{$x$};\draw[->,Slate] (0,-4)--(0,7) node[above]{$y$};
\draw[tmOrange!80!black,line width=0.9pt] (-5,-3) rectangle (10,6);
\draw[tmBlue!80!black,line width=1.2pt] (7,6)--(-5,1)--(-2,-3)--(10,2)--cycle;
\node[above] at (7,6){$A$};\node[left] at (-5,1){$B$};\node[below] at (-2,-3){$C$};\node[right] at (10,2){$D$};
\node[font=\footnotesize] at (-2.5,4.6){$30$};\node[font=\footnotesize] at (-4.2,-2){$6$};
\node[font=\footnotesize] at (6.5,-1.6){$30$};\node[font=\footnotesize] at (9.1,4.9){$6$};
\end{scope}\end{figC}

\step{1}{The enclosing rectangle}
The smallest $x$-coordinate is $-5$, the largest is $10$; the smallest $y$-coordinate is $-3$, the largest is $6$.
\begin{recallbox}The distance between two points on a number line is the larger number minus the smaller: from
$-5$ to $10$ it is $10-(-5)=15$.\end{recallbox}
The rectangle has length $10-(-5)=15$, width $6-(-3)=9$, and area $15\times 9=135$.

\step{2}{The four corner right triangles}
\begin{recallbox}Area of a right triangle $=\tfrac12\times$ the product of the two legs.\end{recallbox}
\begin{itemize}[leftmargin=1.4em,itemsep=1pt]
\item Top-left corner (between $B$ and $A$): legs $7-(-5)=12$ and $6-1=5$: area $\tfrac{12\cdot5}{2}=30$.
\item Bottom-left corner (between $B$ and $C$): legs $-2-(-5)=3$ and $1-(-3)=4$: area $\tfrac{3\cdot4}{2}=6$.
\item Bottom-right corner (between $C$ and $D$): legs $10-(-2)=12$ and $2-(-3)=5$: area $30$.
\item Top-right corner (between $D$ and $A$): legs $10-7=3$ and $6-2=4$: area $6$.
\end{itemize}
Total: $30+6+30+6=72$.

\step{3}{Area of the quadrilateral}
$135-72=63$.

\begin{warnbox}With negative coordinates it is easy to get a length wrong, e.g.\ writing $7-5=2$ instead of
$7-(-5)=12$. Count squares on the picture to check.\end{warnbox}
\verify{In fact $ABCD$ is a parallelogram ($A\to B$ and $D\to C$ are both ``$12$ left, $5$ down''). Triangle
$ABD$ has half the area of the quadrilateral; enclose it in a $15\times 5$ rectangle (from $y=1$ to $6$):
$75-30-\tfrac{15\cdot 1}{2}-6=31.5$. Doubling: $63$. It matches.}
Option \textbf{(B)}.
\ans{(B) $63$}
}

% ---------------------------------------------------------------- 27
\begin{probbox}{27}{\relax}
The figure below is made of $8$ identical small parallelograms joined together. Including these $8$ small
parallelograms, how many parallelograms are there in the figure altogether?
\begin{figC}\begin{scope}[x=0.6cm,y=0.6cm]
\foreach \k in {0,...,4}{\draw[Ink,line width=0.9pt] (2.2*\k,0)--(2.2*\k+2,2);}
\foreach \j in {0,1,2}{\draw[Ink,line width=0.9pt] (\j,\j)--(\j+8.8,\j);}
\end{scope}\end{figC}
\choices{$29$}{$30$}{$26$}{$27$}{$28$}
\end{probbox}
\probsol{
\textbf{Understanding the problem.} Every parallelogram in the figure has two sides on two \emph{horizontal}
lines and two sides on two \emph{slanted} lines of the figure. We count systematically so as not to miss or
repeat any.

\step{1}{Count by height and width (in small cells)}
\emph{Parallelograms $1$ row high} (there are $2$ rows to choose from). A row has $4$ cells side by side:
\begin{itemize}[leftmargin=1.4em,itemsep=1pt]
\item $1$ cell wide: $4$; $2$ cells wide: $3$; $3$ cells wide: $2$; $4$ cells wide: $1$;
\item each row gives $4+3+2+1=10$, and the two rows give $20$.
\end{itemize}
\emph{Parallelograms $2$ rows high}: likewise $4+3+2+1=10$.

Total: $20+10=30$.

\step{2}{A faster count: choose the lines}
\begin{proofbox}A parallelogram is completely determined by choosing $2$ of the $3$ horizontal lines (top and
bottom sides) and $2$ of the $5$ slanted lines (left and right sides). Conversely, each such choice gives exactly
one parallelogram.\end{proofbox}
\begin{recallbox}The number of ways to choose $2$ lines out of $k$ lines is $\dfrac{k\times(k-1)}{2}$ (each pair
is counted twice if we pick a ``first line'' and a ``second line'').\end{recallbox}
Choose $2$ of the $3$ horizontal lines: $\tfrac{3\cdot 2}{2}=3$ ways. Choose $2$ of the $5$ slanted lines:
$\tfrac{5\cdot4}{2}=10$ ways. Number of parallelograms: $3\times 10=30$.

\begin{warnbox}Students often forget the parallelograms $2$ rows high, or the biggest one (the whole
figure).\end{warnbox}
\verify{The two counting methods give the same result, $30$.}
Option \textbf{(B)}.
\ans{(B) $30$}
}

% ---------------------------------------------------------------- 28
\begin{probbox}{28}{\relax}
A right triangular prism has volume $120$ cm$^3$. Its base is a right triangle whose two legs are $3$ cm and
$4$ cm (as shown). What is the height of the prism, in centimetres?
\begin{figC}\begin{scope}[x=0.75cm,y=0.75cm]
\coordinate (A) at (0,2.8);\coordinate (B) at (1.9,3.8);\coordinate (C) at (4.3,2.2);
\coordinate (A1) at (0,0);\coordinate (B1) at (1.9,1.0);\coordinate (C1) at (4.3,-0.6);
\fill[tmBlue!15] (A)--(B)--(C)--cycle;
\draw[Ink,line width=1pt] (A)--(B)--(C)--cycle (A)--(A1)--(C1)--(C);
\draw[Ink,dashed,line width=0.8pt] (B)--(B1) (A1)--(B1)--(C1);
\draw[Ink] ($(B)!0.3cm!(A)$)--($(B)!0.3cm!(A)+(B)!0.3cm!(C)-(B)$)--($(B)!0.3cm!(C)$);
\draw[Ink] ($(B1)!0.3cm!(A1)$)--($(B1)!0.3cm!(A1)+(B1)!0.3cm!(C1)-(B1)$)--($(B1)!0.3cm!(C1)$);
\node[above left] at ($(A)!0.5!(B)$){$3$ cm};\node[above right] at ($(B)!0.5!(C)$){$4$ cm};
\end{scope}\end{figC}
\choices{$12$}{$20$}{$10$}{$16$}{$8$}
\end{probbox}
\probsol{
\textbf{Understanding the problem.} We know the volume and the base (a right triangle). We must find the height of
the prism, i.e.\ the length of a lateral edge (joining the top base to the bottom base).

\begin{recallbox}Volume of a right prism $=$ area of the base $\times$ height:
$V=A_{\text{base}}\cdot h$. Therefore $h=V\div A_{\text{base}}$.\end{recallbox}

\step{1}{Area of the base}
The base is a right triangle; taking one leg as the base, the other leg is the height:
\[ A_{\text{base}}=\frac{3\times 4}{2}=6 \text{ (cm}^2). \]

\step{2}{Height of the prism}
$h=120\div 6=20$ (cm).

\begin{warnbox}Do not confuse ``the height of the prism'' (the distance between the two bases) with a height of
the base triangle. And do not forget to divide by $2$ for the area of the triangle: using $3\times 4=12$ gives $10$
-- the ``trap'' option (C).\end{warnbox}
\verify{$6\times 20=120$ cm$^3$. Correct.}
Option \textbf{(B)}.
\ans{(B) $20$ cm}
}

% ---------------------------------------------------------------- 29
\begin{probbox}{29}{\stars{1}}
Small cubes with edge $1$ are used to build a rectangular box of size $8\times 8\times n$ ($n$ a positive
integer). Let $A$ be the surface area of the box and let $B$ be the total surface area of all the small cubes
(computed for each cube separately and then added up). Find the sum of all values of $n$ for which $\dfrac{B}{A}$
is an integer.
\choices{$86$}{$90$}{$70$}{$78$}{$96$}
\end{probbox}
\probsol{
\textbf{Understanding the problem.} We write $A$ and $B$ in terms of $n$, form the ratio $\dfrac BA$, and find the
values of $n$ that make it an integer.

\step{1}{Compute $A$}
\begin{recallbox}A rectangular box of size $a\times b\times c$ has $6$ faces forming three pairs of equal faces,
so its surface area is $2(ab+bc+ca)$.\end{recallbox}
$A=2(8\cdot 8+8\cdot n+8\cdot n)=2(64+16n)=128+32n$.

\step{2}{Compute $B$}
Number of small cubes: $8\cdot 8\cdot n=64n$. Each cube has $6$ faces of area $1$, so
$B=64n\cdot 6=384n$.

\step{3}{Simplify the ratio}
\[ \frac{B}{A}=\frac{384n}{128+32n}=\frac{32\cdot 12n}{32\cdot(4+n)}=\frac{12n}{n+4}. \]
We ``split off the whole part'': since $12n=12(n+4)-48$,
\[ \frac{12n}{n+4}=\frac{12(n+4)-48}{n+4}=12-\frac{48}{n+4}. \]

\step{4}{When is the ratio an integer?}
$12-\dfrac{48}{n+4}$ is an integer exactly when $\dfrac{48}{n+4}$ is an integer, i.e.\ when $n+4$ is a
\emph{divisor} of $48$.
\begin{recallbox}The positive divisors of $48$: $1,2,3,4,6,8,12,16,24,48$.\end{recallbox}
Since $n\ge 1$, we have $n+4\ge 5$, so we take only the divisors $6,8,12,16,24,48$, which give
\[ n=2,\ 4,\ 8,\ 12,\ 20,\ 44. \]

\step{5}{Add them up}
$2+4+8+12+20+44=90$.

\begin{center}\begin{tabular}{|c|c|c|c|c|c|c|}\hline
$n$ & $2$ & $4$ & $8$ & $12$ & $20$ & $44$\\\hline
$B/A=\dfrac{12n}{n+4}$ & $\tfrac{24}{6}=4$ & $\tfrac{48}{8}=6$ & $\tfrac{96}{12}=8$ & $\tfrac{144}{16}=9$ & $\tfrac{240}{24}=10$ & $\tfrac{528}{48}=11$\\\hline
\end{tabular}\end{center}
\begin{warnbox}Do not forget that $n+4$ can be $48$ (i.e.\ $n=44$), a rather large value. And remember to reject the
divisors $1,2,3,4$, because then $n\le 0$.\end{warnbox}
\verify{The table shows that all $6$ values give an integer ratio. Try another value, $n=1$: $\tfrac{12}{5}$ is not
an integer. Correct.}
Option \textbf{(B)}.
\ans{(B) $90$}
}

% ---------------------------------------------------------------- 30
\begin{probbox}{30}{\relax}
Rectangle $PQRS$ is made up of six identical squares. Points $U,V,W,X,Y,Z$ are midpoints of sides of the squares,
as shown. Which of the following triangles has the greatest area?
\begin{figC}\begin{scope}[x=0.6cm,y=0.6cm]
\draw[Ink,line width=1pt] (0,0) rectangle (6,4);
\draw[Ink,line width=0.9pt] (2,0)--(2,4) (4,0)--(4,4) (0,2)--(6,2);
\foreach \p in {(1,4),(3,4),(4,3),(5,2),(4,1),(6,1)}{\fill[Ink] \p circle (2.2pt);}
\node[above] at (1,4){$V$};\node[above] at (3,4){$W$};\node[right] at (4,3){$X$};
\node[below] at (5,2){$Y$};\node[left] at (4,1){$U$};\node[right] at (6,1){$Z$};
\node[above left] at (0,4){$Q$};\node[above right] at (6,4){$R$};
\node[below left] at (0,0){$P$};\node[below right] at (6,0){$S$};
\end{scope}\end{figC}
\choices{$PVU$}{$PXZ$}{$PVX$}{$PYS$}{$PQW$}
\end{probbox}
\probsol{
\textbf{Understanding the problem.} We need to compute (or compare) the areas of $5$ triangles. To avoid
fractions, let each small square have side $2$; then all the midpoints have whole-number coordinates. Put
$P(0,0)$, $S(6,0)$, $R(6,4)$, $Q(0,4)$; then $V(1,4)$, $W(3,4)$, $X(4,3)$, $Y(5,2)$, $U(4,1)$, $Z(6,1)$.

\begin{recallbox}Area of a triangle $=\tfrac12\times$ base $\times$ height. For a ``slanted'' triangle, enclose it
in a rectangle with sides along the grid and subtract the corner right triangles (as in Problem 26).\end{recallbox}

\step{1}{Two triangles with a side on a side of the rectangle}
\begin{itemize}[leftmargin=1.4em,itemsep=1pt]
\item $PYS$: base $PS=6$, height = distance from $Y$ to $PS$, which is $2$: $A=\tfrac{6\cdot 2}{2}=6$.
\item $PQW$: base $PQ=4$, height = distance from $W$ to $PQ$, which is $3$: $A=\tfrac{4\cdot 3}{2}=6$.
\end{itemize}

\step{2}{Three ``slanted'' triangles}
\begin{itemize}[leftmargin=1.4em,itemsep=1pt]
\item $PVU$: enclosing rectangle $4\times 4=16$; subtract $\tfrac{1\cdot 4}{2}=2$, $\tfrac{3\cdot 3}{2}=4.5$,
$\tfrac{4\cdot 1}{2}=2$. What remains is $16-8.5=7.5$.
\item $PXZ$: enclosing rectangle $6\times 3=18$; subtract $\tfrac{4\cdot3}{2}=6$, $\tfrac{2\cdot2}{2}=2$,
$\tfrac{6\cdot1}{2}=3$. What remains is $18-11=7$.
\item $PVX$: enclosing rectangle $4\times 4=16$; subtract $\tfrac{1\cdot4}{2}=2$, $\tfrac{3\cdot1}{2}=1.5$,
$\tfrac{4\cdot3}{2}=6$. What remains is $16-9.5=6.5$.
\end{itemize}
\begin{figC}\begin{scope}[x=0.6cm,y=0.6cm]
\fill[tmOrange!15] (0,0) rectangle (4,4);
\fill[tmBlue!35] (0,0)--(1,4)--(4,1)--cycle;
\draw[Line] (0,0) grid[step=1] (6,4);
\draw[Ink,line width=1pt] (0,0) rectangle (6,4);
\draw[tmBlue!80!black,line width=1.1pt] (0,0)--(1,4)--(4,1)--cycle;
\node[font=\footnotesize] at (0.35,3.4){$2$};\node[font=\footnotesize] at (3,3){$4.5$};\node[font=\footnotesize] at (3,0.35){$2$};
\node[below left] at (0,0){$P$};\node[above] at (1,4){$V$};\node[right] at (4,1){$U$};
\end{scope}\end{figC}

\step{3}{Compare}
$7.5>7>6.5>6=6$. The largest triangle is $PVU$.

\begin{warnbox}Do not guess by eye: triangle $PXZ$ looks ``big'' because it stretches all the way to side $RS$, but
its area is only $7$, less than that of $PVU$.\end{warnbox}
\verify{If the side of each square is $1$, all areas are divided by $4$ ($1.875$; $1.75$; $1.625$; $1.5$; $1.5$),
and the order does not change.}
Option \textbf{(A)}.
\ans{(A) triangle $PVU$}
}
